\clearpage\begingroup\SourceLeaf{106b}\HistoricalHead{101}{R}{GAUSS AN SCHUMACHER UND GERLING.}
\begin{center}[18.]\end{center}
\begin{center}\textsc{Gauss} an \textsc{Gerling}.\quad Göttingen, 26. August 1835.\end{center}
\noindent .\ .\ .\ .\ .\ .\ .\ .\ .\ .\ .\ .\ .\ .\ .\ .\ .\ .\ .\ .\ .\par
\indent In Beziehung auf die magnetogalvanische Telegraphie bin ich vor einigen
Wochen auf eine ganz neue Art die Induction dazu zu verwenden gekommen,
\clearpage\SourceLeaf{107a}\HistoricalHead{102}{L}{PHYSIK. BRIEFWECHSEL.}
\noindent die sich in der Anwendung ungemein zierlich macht. Es lassen sich, wie die
Erfahrung gelehrt hat, in der Minute 6 bis 8 Buchstaben transmittiren. Sollte
jedoch dergleichen Telegraphie für bürgerliche Zwecke im Ernst und als bestehende
Staatseinrichtung benutzt werden, so würde es vortheilhafter sein,
anstatt \textls[120]{zweier} Stränge (Einer Kette) ihrer mehrere z. B. 7 zwischen den
beiden Orten gehen zu lassen. Freilich werden dann die Kosten für Metall
in demselben Verhältniss höher, dagegen kann man dann aber die Einrichtungen
leicht so machen, dass erstlich die Transmission noch viel geschwinder vor
sich geht, und zweitens, was noch wichtiger sein möchte, dass die currente
Handhabung durch ganz ordinaire Personen geschehen kann. Ich bin überzeugt,
dass auf solche Weise die elektromagnetische Telegraphie theils vollkommen
eben so schnell wie die optische, wo nicht schneller, ausgeführt
werden kann, theils mit \textls[120]{viel} geringeren Kosten, da man bei jener ohne Frage
Zwischenräume von 20 und mehreren Meilen nehmen kann, und drittens mit
anderen Vortheilen, deren man bei der optischen entbehrt, z. B. der Unabhängigkeit
von Wetter und Tageszeit, dem völligen Unbemerktbleiben, dass
correspondirt wird etc. \mbox{.\ .\ .\ .\ .\ .\ .\ .\ .\ .\ .\ .\ .\ .\ .\ .\ .\ .\ .\ .\ .}

Da übrigens bei derartigen Ausführungen die Wissenschaft eigentlich gar
nicht interessirt ist, so scheint mir in der Ordnung, dass jene, obwohl bona
officia zu leisten bereit, doch solche nicht aufdrängt, sondern wartet, bis solche
gesucht werden. \mbox{.\ .\ .\ .\ .\ .\ .\ .\ .\ .\ .\ .\ .\ .\ .\ .\ .\ .\ .\ .\ .}
\clearpage\endgroup
